Our second main result is a description of the rational homotopy type of (a given component of) the space of almost complex structures on six-manifolds, under the additional assumption that $b_1(M) = 0$ and $\int_M c_1c_2 - c_3 \neq 0$ (Theorem~\ref{Qhomotopy6}). For this we use an expression of $J(M)$ as the quotient of the space of sections of an $S^7$-bundle by an $S^1$-action (Theorem~\ref{fibration}), which one can compare with the result obtained for $c_1=0$ in~\cite[Proposition~5.1]{FGM21}.
Throughout we provide examples for all of the mentioned results.

\subsection*{Notation and conventions}

 The projectivization of a complex vector bundle $E$ is denoted by $\mathbb{P}(E)$. We denote the Chern classes of a complex vector bundle by~$c_i$, and the Pontryagin classes of a real vector bundle by~$p_i$; if the vector bundle is complex, by its Pontryagin classes we mean those of the underlying real bundle. The Euler class is denoted by~$\mathsf{e}$. The Euler characteristic of a space is denoted by~$\chi$, and~$\sigma$ is the signature of an oriented closed manifold. All manifolds are assumed connected unless otherwise stated.

We denote by $\Gamma(E)$ the space of sections of a bundle $E \to M$ (where $E$ is not necessarily a~vector bundle). If $s \colon M \to E$ is a section, $\Gamma(E)_s$ denotes the connected component of~$s$ in~$\Gamma(E)$.
For $E \to M$ a vector bundle with a metric, we write $S(E) \to M$ for the associated sphere bundle. We denote by $J(M)$ the connected component of a specific point in $\Gamma(Z_+(M))$, which will be clear from context. For a vector space~$V$ with an inner product, $J(V)$ will denote the space of all orthogonal (linear) complex structures on $V$. If $V$ is oriented, $J_+(V)$ (respectively~$J_-(V)$) denotes the connected component of~$J(V)$ consisting of elements which induce the given orientation on $V$ (respectively the opposite orientation). By $\Map(X,Y)$ we denote the space of unbased maps $X \to Y$ with the compact-open topology.

Cohomology is singular cohomology with integral coefficients unless another choice of coefficients is indicated. The evaluation of a cohomology class $\alpha$ on a homology class $A$ is denoted by $\langle \alpha, A \rangle$ or by $\int_A \alpha$. In the context of rational homotopy theoretic models, $\Lambda(x_i)$ denotes the free graded-commutative algebra on generators $x_i$.

We will make frequent use of the projective bundle formula, which states that for a complex vector bundle $E \to M$ of complex rank $k$, the cohomology of the projectivized bundle $\mathbb{P}(E) \xrightarrow{p} M$ is the free $H^*(M)$-algebra on one generator $x$ subject only to the relation $x^k + c_1(E)x^{k-1} + \cdots + c_{k-1}(E)x + c_k(E) = 0$, where $x$ is the first Chern class of the line bundle $\mathcal{O}_{E^*}(1)$ (see~\cite[Definition~15.13]{D12}). This line bundle is dual
to the tautological (Hopf) line subbundle of $p^*(E)$ whose fiber over $\ell \in \mathbb{P}(E_x)$ is $\ell \subset E_x$. It restricts to $\mathcal{O}(1)$ on each fiber $\CP^{k-1}$.

For a real vector space $V$ with an endomorphism $J$ squaring to $-\mathrm{Id}$, we denote by $\Lambda_\C^*(V)$ the (complex) exterior algebra of the complex vector space $(V,J)$. For an almost complex manifold $(M,J)$, $\Lambda_\C^\ast (TM)$ is the (complex) exterior bundle associated to $(TM,J)$.

